\section{From bounded cohomology to comparison degree}\label{sec:transfer}

The entropy construction supplies a root probability map \(F\) whose
energy \(e_h(F)\) has a large coefficient in the weighted scalar-curvature
inequality.  Its coordinatewise square takes values in the probability
simplex.  We now turn its pairing with the bounded cocycle into comparison
degrees to which Theorem~\ref{thm:degree} applies.  A common integer
clears the rounded cubical coefficients, giving integer multiplicities;
the resulting measured degree may still be real.  Its expectation will
retain the original cocycle pairing exactly.  The different lengths
assigned to the tangent and normal directions of the realization are
what remove the auxiliary dimension from the final estimate.

The probability-map and auxiliary-degree strategy is adapted from
Section~3 of \cite{OpenAIScalarInessentiality2026}: Lemmas~3.2 and 3.4
develop the probability map and dual chains, and Proposition~3.6 provides
their pairing.  Lemma~3.8 supplies a transverse parameter cycle with
nonzero intersection count; the proof of Theorem~3.1 constructs the
comparison map and identifies its local degree.  We prove the
quantitative measured transfer and its estimates here.

\begin{theorem}\label{thm:scalar-volume}
For every integer \(n\geq 3\), there is a finite constant \(C_n>0\) such
that every closed connected oriented smooth \(n\)-manifold \(M\) and every
smooth Riemannian metric \(g\) with \(\Scal_g\equiv-1\) satisfy
\[
                         \|M\|\leq C_n\Vol_g(M).
\]
\end{theorem}

We prepare the proof for a manifold with \(\|M\|>0\).  Choose the
homogeneous alternating bounded cocycle \(c\), the initial probability
map, and the number \(I\geq\|M\|/2\) provided by
Proposition~\ref{prop:detector}; in particular \(\|c\|_\infty\leq1\).
We first consider conditional data having the properties asserted by
Theorem~\ref{thm:entropy}, without yet invoking their existence.
Denote the scalar constant in this contract by \(C_{\rm ent}\).
All functions called elementary in this section have a fixed height of
exponentiation on a polynomial, with parameters depending only on \(n\).

Write \(\mathcal M=(\widetilde M\times\Omega)/\Gamma\) for such a measured
family.  Denote its metric, weight, density, and
root probability map by \(h,\phi,\rho,F\), and put
\[
 \mathfrak e=e_h(F),\qquad W=\log(dV_h/\rho).
\]
Its estimates are
\begin{equation}\label{eq:transfer-entropy}
 \rho\leq dV_g,\qquad W\geq0,\qquad
 \wscal(h,\phi)\geq-C_{\rm ent}+2W-n\log E+E\mathfrak e,
 \qquad \mathfrak e^{1/2}\leq E^{-1/100}.
\end{equation}
There is a fixed invariant graph on the countable label set, with finitely
many free \(\Gamma\)-orbits, containing every cooccurrence of nonzero
coordinates of \(F\).  The integer \(q\geq n+1\) bounds the number of
orbits and the maximum graph degree plus one.  The equivariant coloring
by \(q\) colors separates any two labels at graph distance at most two.
Its dependence on \(\omega\in\Omega\) is Borel.  Set
\(m=q-n>0\).  Integrals over \(\mathcal M\) and over the measured spaces
constructed below include their transverse measures.

\medskip\noindent\emph{Conditional transfer target.}
For every \(n\geq3\), there are a constant \(K_n>0\) and an elementary
function \(\Psi_n:[1,\infty)\to(0,\infty)\), chosen from \(n\) alone,
such that every choice of detector and conditional entropy data just
specified, including its constant \(C_{\rm ent}\), satisfies
\begin{equation}\label{eq:transfer-conditional-target}
 E\geq\Psi_n(q)\quad\Longrightarrow\quad
       I\leq K_n e^{C_{\rm ent}/2}\int_{\mathcal M}\rho.
\end{equation}
Thus \(K_n\) and \(\Psi_n\) are fixed before the manifold, cocycle,
instance geometry, and \(C_{\rm ent}\) are chosen.  We carry out the
estimates for each integer \(q\geq n+1\), uniformly over conditional
data satisfying the stated label and coloring bounds.
The grid, the common rounding, and an upper density bound for the
resolution come first.  That density bound permits a small collar;
the ensuing curvature and rotation bounds then determine one comparison
width and the required energy threshold.  Only after the conditional
target is closed will Theorem~\ref{thm:entropy} supply an actual pair
\((q,E)\) and make \(C_{\rm ent}\) dimension-only.

\subsection{A flat neighborhood and a cubical cocycle}

For a fixed \(\omega\), let \(K=K_\omega\) consist of all nonnegative
unit root vectors whose support is a clique in the label graph.  Thus
\(K\) is the union of the spherical orthants on those cliques.  A vector
\(y\in K\) has a color image \(z_y\in\mathbb R^q\), obtained by placing
each of its root coordinates at the color of its label.  Distinct labels
in a clique have distinct colors.  Distant cliques can reuse colors, so
the flat neighborhood below keeps distinct label configurations with the
same color image on separate sheets for comparison, while retaining
Euclidean coordinates for the grid averages.

\begin{lemma}[Flat sheet neighborhood]\label{lem:transfer-flat}
Put \(\delta=(1000\sqrt q)^{-1}\).  There is an oriented flat smooth
\(q\)-manifold \(\mathcal U_\omega\) with a local isometry
\(z:\mathcal U_\omega\to\mathbb R^q\), and an inclusion of \(K_\omega\)
into it, such that the image of every \(y\in K_\omega\) has a flat ball
chart of radius \(10\delta\) centered at \(z_y\).  This construction is
a standard Borel family in \(\omega\), with countably many charts and an
equivariant \(\Gamma\)-action.  The map \(F\) lifts smoothly to this
family, and, writing \(z_F=z\circ F\),
\[
                         e_h(z_F)=e_h(F)=\mathfrak e.
\]
There is a closed smooth \(n\)-form, again denoted \(\eta\), obtained
in these flat charts by pulling back the detector form through the cutoff
probabilities constructed below.  It is Borel and equivariant in the
family, and satisfies
\begin{equation}\label{eq:transfer-form}
        \int_{\mathcal M}F^*\eta=I,\qquad |\eta|_{\HS}\leq C_n.
\end{equation}
Here the exterior Hilbert norm uses the Euclidean metric on the charts.
For every fixed derivative order its derivatives admit an elementary
bound in \(q\).
\end{lemma}

\begin{proof}
Take one copy of \(B(z_y,10\delta)\subset\mathbb R^q\) for each witness
\(y\in K\).  At the same point \(z\) in two such balls, identify the two
copies if their witnesses share a label in their nonzero supports.  We
check that these identifications give ball charts, rather than a branched
space.  For every point under consideration, some coordinate of \(z\)
is at least \(q^{-1/2}-10\delta\).  Every witness for that point has a
nonzero label at this color.  If two witnesses share any label, their
labels at every common color agree: each of the two labels at that color
is adjacent to the shared label, so distance-two separation of colors
forces them to be the same.  Conversely, agreement at the indicated
large color is already a shared label.  Thus, at \(z\), the proposed
identification is precisely agreement on the label at that large color,
which is an equivalence relation.

The same color remains positive in all witnesses at every sufficiently
nearby point; indeed the margin \(q^{-1/2}-21\delta>0\) suffices after
a further displacement of at most \(\delta\).  In that neighborhood
the equivalence classes are still distinguished by its label.  Each ball
therefore projects injectively as a local chart.  Points with distinct
Euclidean projections can be separated in Euclidean space, and points
with the same projection and different labels at the large color have
disjoint chart neighborhoods just described.  This proves the Hausdorff
property.  Rational dense choices of centers in each spherical orthant
give a countable subatlas: there are countably many labels and cliques.
The descriptions by labels, clique coordinates, and color assignments
are Borel, and are unchanged by simultaneous translation of labels and
\(\omega\).  They prove all the claimed family and equivariance
properties.  At points of \(K\) this also shows that its inclusion is
unambiguous.

The lift of \(F\) is locally in one of these charts: its label at a
chosen large color persists in a neighborhood.  A nonnegative smooth
root coordinate has zero differential wherever it is zero.  All the
nonzero coordinates at a point have different colors.  It follows that
the nonzero rows of \(dF\) are merely moved to different color positions
by \(dz_F\), proving the asserted energy equality.

To construct this cutoff pullback, choose a nonnegative smooth function \(\chi\)
which is zero for \(t\leq12\delta\), equals \(t\) for
\(t\geq25\delta\), lies between zero and \(\max\{t,0\}\), and satisfies the usual rescaled derivative bounds
\(|\chi^{(j)}|\leq C_j\delta^{1-j}\) for \(j\geq1\).  In a ball put
\[
       w(z)=\frac{(\chi(z_j))_{j=1}^q}{|(\chi(z_j))_{j=1}^q|},
       \qquad p_j=w_j^2.
\]
The denominator is bounded below by an absolute positive constant.
In fact \(z\) is within \(10\delta\) of a nonnegative unit vector, and
the change in each coordinate caused by cutting it off is at most
\(25\delta\); hence its total change is at most \(25\delta\sqrt q\).
If \(\chi(z_j)\ne0\), every witness in this ball has a positive
coordinate at color \(j\), because its distance from \(z\) is less
than \(10\delta\).  The label for every active coordinate is therefore
unambiguous on chart overlaps.  Evaluate
\[
                         \eta=n!\,c(p,dp,\ldots,dp)
\]
using those labels.  The cocycle law makes this form closed, as in
Proposition~\ref{prop:detector}, and a straight probability homotopy
from the old probabilities to these new ones preserves the measured
pairing.  Its local finiteness and smooth bounds are uniform for the
instance, so the measured Stokes argument of that proposition applies.
This gives the first identity in \eqref{eq:transfer-form}.

For the dimension-only norm bound, first regard the form on nonnegative
unit root vectors \(w\).  In the Euclidean coordinate directions
\(i_1,\ldots,i_n\), its component is
\[
 n!2^n\Bigl(\prod_{j=1}^n w_{i_j}\Bigr)
                \sum_a w_a^2 c(a,i_1,\ldots,i_n).
\]
Its absolute value is at most \(n!2^n\prod_j w_{i_j}\).  The sum of the
squares of these products over all increasing \(n\)-tuples is at most
\((\sum_i w_i^2)^n=1\).  Thus its exterior Hilbert norm is at most a
constant depending only on \(n\).  The map \(z\mapsto w(z)\) has
absolutely bounded operator norm of its differential: its denominator
is bounded below, normalization is an orthogonal projection divided by
that denominator, and \(\chi'\) is absolutely bounded.  Pullback
multiplies the exterior Hilbert norm by at most the \(n\)-th power of
this operator bound.  This proves the second assertion in
\eqref{eq:transfer-form}.  Differentiating the same formulas a fixed
number of times gives elementary bounds (polynomial bounds in \(q\)
already suffice), since \(\delta^{-1}=1000\sqrt q\).
\end{proof}

Choose a rotation \(V\in SO(q)\) with Haar probability, independently
of the entropy data, and put \(\widehat z=Vz\).  In these coordinates
take the cubical grid of step \(r\) and shift \(s\), where \(s\) is
uniform in \([0,r)^q\).  Fix \(r=r(q)>0\) so that
\(r\sqrt q\ll\delta\) and so that every \(n\)-component of \(\eta\)
in every rotated frame changes by at most \(q^{-n}\) over a ball of
radius \(10r\sqrt q\) in the inner sheet charts.  The derivative bounds
of Lemma~\ref{lem:transfer-flat} permit such a choice with an elementary
bound on \(r^{-1}\), uniformly in \(V,s,\omega\).

A primal \(n\)-cell has an axis set \(\mathcal I\subset\{1,\ldots,q\}\)
of size \(n\), varies along those grid intervals, and has the remaining
coordinates at grid vertices.  Orient it in increasing axis order.
Include each lift \(C\) whose center is at distance less than \(4\delta\)
from some witness \(z_y\) in its sheet chart.  The full cell fits in that
chart.  Its dual \(m\)-cube \(C^\#\) has the same center, varies along
the complementary axes by \([-r/2,r/2]\), and has the other coordinates
fixed.  It too fits in the chart.  Orient it so that the primal oriented
basis, followed by the negatives of its dual oriented basis, is a
positive basis of \(\mathbb R^q\).  Give the primal cell the coefficient
\begin{equation}\label{eq:transfer-cell-value}
                              u_C=\int_C\eta.
\end{equation}

These coefficients satisfy the cubical cocycle equations where they
are needed.  If a primal \((n+1)\)-cell has its center within
\(2\delta\) of a witness, all its boundary \(n\)-faces are included
with coherent lifts, and their signed integrals sum to zero by Stokes.
A facet of a dual cube is centered at just such a primal
\((n+1)\)-cell.  The incidence signs of its incident dual cubes agree,
up to one common sign, with the boundary signs of the primal cell:
this follows by placing the varying primal axes first and the negative
dual axes last in the ambient oriented basis.  Consequently the dual
\(m\)-cubes weighted by \(u_C\) have balanced oriented boundary near
\(K\), in particular at every facet meeting the tube of radius
\(\delta\).  Indeed the corresponding centers differ from those
facet points by at most \(r\sqrt q\), so they have the stipulated
\(2\delta\) witnesses.  Projected grid cells on other sheets give no
incidence: all these statements use lifts in the same flat chart.

\begin{lemma}[Finite-type unbiased rounding]\label{lem:transfer-rounding}
There is an integer \(D_0\geq1\), bounded by an elementary function of
\(q\) and constant in \(V,s,\omega\), and a Borel equivariant random
choice of coefficients \(u'_C\in D_0^{-1}\mathbb Z\) such that
\begin{equation}\label{eq:transfer-rounding}
 \E_{\rm rd}u'_C=u_C,\qquad |u'_C-u_C|\leq r^nq^{-n}.
\end{equation}
The expectation is conditional on the grid and entropy data.  Every
cocycle equation just described is satisfied exactly by the \(u'_C\).
\end{lemma}

\begin{proof}
The closure equations couple faces whose anchors can differ.  We encode
all local coefficients in one ordered finite vector and round within the
joint kernel of its incidence equations, so every cell reads from the
same rounded vector.  We first establish that the vector has elementary
size despite the infinitely many lifted cells.

All projected centers lie
in a fixed bounded Euclidean ball.  For fixed \(V,s\), their grid
indices and axis sets therefore have at most \((C/r)^{Cq}\) choices.
At a selected center choose, breaking ties by color order, its largest
unrotated color coordinate.  This coordinate is at least
\(q^{-1/2}-4\delta\), and its label occurs in every witness for the
center.  It determines the lift uniquely, by the proof of
Lemma~\ref{lem:transfer-flat}; call it the anchor label.

The closed graph one-neighborhood of an anchor has at most \(q\)
labels.  Every center-test witness, meaning a witness at distance less
than \(4\delta\) from the center, lies in this neighborhood because it
contains the anchor.  Every point of the full cell is within
\(4\delta+r\sqrt n/2<5\delta\) of each such witness.  Thus every
active color, whose coordinate is greater than \(12\delta\), occurs
in each center-test witness.  For an arbitrary ball-chart witness of
the center the corresponding bound is instead
\(10\delta+r\sqrt n/2<11\delta\), which is still below the active
cutoff and gives the same conclusion.  Thus the
coefficient \(u_C\) is determined by its projected grid index, its
anchor, and the color assignment on this one-neighborhood.  In
particular, no labels farther out in the graph enter its calculation.

Choose and order representatives of the at most \(q\) anchor orbits,
and order each of their finite one-neighborhoods.  An anchor has a
unique translate to its representative because the label action is
free.  After this translate, its neighborhood is a fixed finite set
and the cocycle evaluations are unchanged by homogeneity.  Fix one
ordered list of slots consisting of an anchor representative, an
injective color assignment on its neighborhood, an axis set and grid
index, and a witness clique in that neighborhood containing the anchor.
It is enough to use the fixed index box
\(|k_j|\leq\lceil3/r\rceil+3\), which contains every possible selected
center for all shifts.  The slot is valid when the anchor color is the
canonically chosen largest unrotated color of its center and that center
is at distance less than \(4\delta\) from the spherical orthant of its
witness clique.  Give a valid slot the integral of the explicit form
over its primal cell, using the labels specified by the slot, and give
an invalid slot value zero.  The preceding active-color observation
makes this integral unambiguous.  The validity tests and integrals are
Borel in \(V,s\).

This is one vector \(v=v(V,s)\in\mathbb R^{N_0}\) for the fixed graph
and cocycle, common to all \(\omega\).  Only its size and the bounds
below are uniform over all graphs and cocycles of the given complexity:
there are at most \(q^q\) neighborhood color assignments, \(2^q\)
witness cliques, and \((C/r)^{Cq}\) projected indices, so \(N_0\) is
elementarily bounded in \(q\).  An actual cell receives the slot
\(t(C,\omega)\) obtained by normalizing its canonical anchor, restricting
its actual coloring, and choosing the first valid witness clique in the
fixed order.  Its value is \(u_C=v_{t(C,\omega)}\).  This type is Borel
and is unchanged by translating \(C\) and \(\omega\) together.  Slots
with the same anchor, restricted coloring, axes, and grid index but
different valid witness cliques have the same value: the
\(5\delta\) observation gives the identical active-label formula
throughout the primal cell.

Impose on this list every integer row \(b\) with
\(\|b\|_\infty\leq2(n+1)\) for which \(b\cdot v=0\).  There are
finitely many candidate rows, so the resulting matrix depends Borelly
on \(v\).  In particular it includes the difference row for any two
duplicate valid slots.  It also includes every needed cubical equation,
even when its faces have anchors normalized by different deck elements:
all their values are entries of this same vector \(v\), and Stokes gives
the signed zero sum of those entries.  After repeated slots in its
\(2(n+1)\) faces have been combined, each coefficient has the stated
bound.  The kernel of this
integer matrix has a real basis of integer vectors \(b_1,\ldots,b_k\)
with
\[
                   k\leq N_0,\qquad
        \max_j\|b_j\|_\infty\leq K_0:=(C_nN_0)^{C N_0}.
\]
To see this, choose a nonsingular maximal minor and solve for its pivot
coordinates in terms of the free coordinates; multiplying each basis
vector by its nonzero determinant makes it integral.  Hadamard's
bound for minors of a matrix with entries at most \(2(n+1)\) gives
the displayed estimate.  Lexicographically choosing the pivot minor
makes the basis Borel.  The case of a zero-dimensional kernel simply
has \(v=0\) and needs no rounding.

Write \(v=\sum_j t_jb_j\), using the pivot solution, which is Borel.
Independently round each \(t_j\) to the two adjacent multiples of
\(D_0^{-1}\), choosing their probabilities so that its conditional
expectation is \(t_j\).  Choose a single integer
\(D_0\geq N_0K_0r^{-n}q^n\), using uniform elementary upper bounds
for \(N_0,K_0\).  The resulting vector \(v'\) lies in the same
kernel, has entries in \(D_0^{-1}\mathbb Z\), has expectation \(v\),
and satisfies \(\|v'-v\|_\infty\leq N_0K_0/D_0\leq r^nq^{-n}\).
All operations, including taking floors and using independent uniform
parameters for the adjacent choices, are Borel.  The fixed slot order
is used for every pivot choice.  Thus one parameter \(\xi\), independent
of \(\Omega\), gives one rounded vector \(v'(V,s,\xi)\) for all cells.
Set \(u'_C=v'_{t(C,\omega)}\).  The invariance of the type under
translation proves equivariance, and the common zero rows prove all
the asserted incidence equations.
\end{proof}

The rounded coefficients define the original oriented dual current
\begin{equation}\label{eq:transfer-rounded-current}
                    Z_{\rm rd}:=\sum_C D_0u'_C[C^\#].
\end{equation}
Here \([C^\#]\) is integration over the oriented cube, and the
coefficients are integers.  We verify the local finiteness implicit
in this notation.  The form bound and rounding give
\(|u_C|\leq C_nr^n\) and an elementary bound for each \(|D_0u'_C|\).
Let \(B\) be a small flat ball in one sheet chart and \(Q\) one full
projected dual cube.  The intersection \(Q\cap z(B)\) is convex.  Any
lift of \(Q\) meeting \(B\) agrees with the chart inverse on this
connected intersection, by uniqueness of lifts for the local
diffeomorphism \(z\).  Thus two lifts meeting \(B\), even at different
projected points, agree there and then agree on the connected cube
\(Q\).  At most one geometric lift of each projected cube meets \(B\).
Only finitely many projected grid cubes meet \(z(B)\), with bounded
coefficients, proving neighborhood local finiteness.

Let \(\mathcal T_\delta\subset\mathcal U_\omega\) be the strict tube
of points at distance less than \(\delta\) from a witness in its ball
chart.  The exact incidence equations in
Lemma~\ref{lem:transfer-rounding} give
\begin{equation}\label{eq:transfer-current-closed}
                    \partial Z_{\rm rd}\big|_{\mathcal T_\delta}=0.
\end{equation}
This is a local assertion: the boundary current vanishes on test forms
compactly supported in \(\mathcal T_\delta\).

To resolve \(Z_{\rm rd}\) smoothly, realize each coefficient by
\(|D_0u'_C|\) oriented copies of its cube, reversing orientation when
\(u'_C<0\).  At every facet where boundary balance is required, pair
positive and negative incident copies and glue them by their common
face.  Leave other facets free, and double this diagram with an
oppositely oriented mirror, joining each free facet to its mirror.
All codimension-one faces of the doubled diagram are now paired.
Local axis, grid, and copy indices order the incident copies, so these
choices may be Borel and equivariant.

The doubling does not join the two sides near \(K\), even through a
lower face.  If a point of any face maps into \(\mathcal T_\delta\),
every path of facet identifications relating chambers incident to that
face preserves its image point.  Each facet on the path meets the tube,
so was balanced before doubling and was not a free mirror facet.  The
entire incidence class at that point stays on its original side or its
mirror side.  At a point of a lifted grid face, only boundedly many
local grid indices and their boundedly many copies can meet.  Facet
identifications preserve that point and its sheet, so even after
doubling they relate only this bounded collection of incident chambers.

Take barycentric subdivision of the doubled cubes.  Each top simplex
\(\sigma\) has vertices colored \(0,\ldots,m\) by the dimensions of
its barycentric faces, and every facet pairing respects these colors.
Write \(\varepsilon_\sigma\in\{1,-1\}\) for its sign relative to that
color order, and call this doubled signed simplicial diagram
\(\mathcal P\).  It maps facewise to \(\mathcal U_\omega\), with
elementary subdivision and incidence bounds.  The original current
\(Z_{\rm rd}\) remains the object to be represented on
\(\mathcal T_\delta\); the doubling supplies closed facet pairings
for the resolution.

\subsection{A controlled smooth measured resolution}

The diagram need not be a manifold at faces of codimension greater
than one.  The following construction replaces it by a measured family
of smooth manifolds while preserving the represented chain.  The
small portion on which its metric is not the pulled-back flat metric
will be controlled by its coordinate volume.

The reflected gluing below uses Gaifullin's explicit permutohedral
resolution, including the face-preserving sign-changing pairings,
their conjugation under reflections, and the auxiliary bit toggles
\cite[pp.~562--563]{Gaifullin2008}.  We replace his rank-indexed bits
by independent subset-indexed bits, put invariant probability laws on
the pairing data, and supply smooth measured charts and quantitative
metric bounds.

\begin{lemma}[Measured resolution]\label{lem:transfer-resolution}
For the diagram \(\mathcal P\) there is a standard Borel measured
family \(\mathcal A\) of oriented smooth \(m\)-manifolds and a smooth
map \(X:\mathcal A\to\mathcal U_\omega\), equivariant and Borel in
all the data.  The integration identities below are for fixed grid,
entropy, and rounding data, including \((V,s,\omega,\xi)\); their
integrals use the conditional state measure and depend Borelly on those
data.  The points over \(\mathcal T_\delta\) represented by
original chamber tiles form a Borel leaf-open subset
\(\mathcal A_{\rm orig}\subset X^{-1}(\mathcal T_\delta)\).  For every
compactly supported smooth \(m\)-form \(\alpha\) on \(\mathcal T_\delta\),
\begin{equation}\label{eq:transfer-local-push}
                 \int_{\mathcal A_{\rm orig}}X^*\alpha=Z_{\rm rd}(\alpha).
\end{equation}
For the full doubled diagram, measured oriented integration also gives
\begin{equation}\label{eq:transfer-push}
                 X_*[\mathcal A]=[\mathcal P].
\end{equation}
For a parameter \(0<\lambda<1/2\) and a choice
\begin{equation}\label{eq:transfer-gamma}
  0<2\gamma_0<(m+1)^{-(m+1)}\lambda^{m(m+1)},
\end{equation}
there are a smooth metric \(k\) and a closed seam set in each tile such
that \(\|dX\|_k\leq1\), and \(k=X^*k_{\rm flat}\) off the seam.
There are buffered smooth charts with elementary overlap, metric,
inverse-metric, and fixed-order derivative bounds whenever
\(\lambda^{-1}\) and \(\gamma_0^{-1}\) are elementary.  In the sorted
charts for each tile, the seam is contained in the union of the slabs
\(|Z_j|\leq\lambda\), \(1\leq j\leq m\).  The \(k\)-volume density
in those charts has an elementary upper bound chosen before
\(\lambda,\gamma_0\); in particular this bound does not depend on
these two parameters.

One may prescribe for each chamber over a dual cube a rotation
\(O_\sigma\in SO(q)\).  There is also a smooth Borel equivariant
\(\mathcal R:\mathcal A\to SO(q)\) equal to \(O_\sigma\) off the
seam in that tile, equal to the identity near all tile facets, and
homotopic to the identity.  Its fixed-order chart derivatives and
\(\|d\mathcal R\|_k\) are elementarily bounded under the same
hypotheses.
\end{lemma}

\begin{proof}
\emph{Invariant reflection states.}
For each nonempty proper subset
\(S\subset\{0,\ldots,m\}\) and each face with exactly those vertex
colors, pair the incident positive chambers with the incident negative
chambers.  The two counts agree.  Indeed choose a color omitted by
\(S\); the given codimension-one matching in \(\mathcal P\) across
facets omitting this color pairs the chambers at that face with opposite
signs.  All incident sets are
finite by the incidence bound above.  Choose uniformly among their
possible pairings, independently for each face and each \(S\), and
let \(p_S\) denote the resulting involution on chambers.  Use also
independent uniform bits \(t_S\in\{0,1\}\), one for each \(S\).
The state space consists of \(\zeta=(\sigma,p,t)\), with counting measure
on \(\sigma\) and these product probabilities on \(p,t\).  These
are countable products of finite probability spaces.  As the original
data vary, the incidence lists and finite pairing choices are Borel,
so the product laws are Borel probability kernels.

Define a move \(s_S\) by
\[
 \sigma\longmapsto p_S\sigma,\qquad t_S\longmapsto1-t_S,
 \qquad p_R\longmapsto p_Sp_Rp_S\quad(R\subsetneq S),
\]
leaving the other data unchanged.  The conjugated \(p_R\) is again
a pairing at each \(R\)-face: \(p_S\) preserves every such face and
flips the chamber sign, so conjugation preserves its face and its
sign-reversing property.  Each \(s_S\) is an involution.  If
\(R\subsetneq S\), the two moves commute.  On chambers, for example,
their two orders both give \(p_Sp_R\sigma\), since after \(s_S\) the
\(R\)-pairing is \(p_Sp_Rp_S\).  On a pairing indexed by a subset of
\(R\), both orders conjugate by \(p_Sp_R\); on the other indices the
same assertion follows directly from which of the two moves changes
that index.  The bits commute as well.  Thus all moves in a nested
flag commute, and they act freely: each changes its own bit, so no
nonempty product of these moves can fix a state.

These moves preserve the state measure.  They permute the counting
chamber index, flip a uniform bit, and, conditional on \(p_S\),
conjugate each relevant uniformly distributed finite pairing by a
permutation of its incident set.  The latter operation permutes its
finite set of possible pairings and preserves its uniform law.  It
also preserves the independent product of those laws.  All moves are
Borel and commute with translating the original label data.

\medskip\noindent\emph{Smooth leaves and the measured atlas.}
Attach the following tile to each state.  Start with the positive
simplex \(x_0,\ldots,x_m>0\), \(\sum_i x_i=1\), and take the closure
of its collection of pair ratios \((x_i/x_j)_{i\ne j}\) in the
product of copies of \([0,\infty]\).  For each permutation
\(i_0,\ldots,i_m\) use the coordinate chart
\begin{equation}\label{eq:transfer-ratios}
                u_j=x_{i_j}/x_{i_{j-1}}\in[0,4),\qquad 1\leq j\leq m.
\end{equation}
All pair ratios in the chart are products of these or their reciprocals,
so the \(u_j\) are independent coordinates.  Sorted orders have
\(u_j\leq1\) and cover the closure.  The facet \(u_j=0\) has label
\(S_j=\{i_0,\ldots,i_{j-1}\}\); the labels of a corner are nested.
Glue each facet to the identical resolved coordinates in the state
obtained by \(s_{S_j}\).  The commutation just proved makes the gluing
consistent at corners.

Use signed square root coordinates \(u_j=Z_j^2\), with
\(Z\in(-2,2)^m\), starting from any state.  For \(Z_j<0\) use the
tile obtained by \(s_{S_j}\); use the product of these commuting moves
when several coordinates are negative.  At zero use the starting tile
as a convention for representing the same glued point.

Figure~\ref{fig:resolution} illustrates the tile and the four signed
quadrants at a two-facet corner when \(m=2\).

\par\medskip
\noindent\begin{minipage}{\textwidth}
\makeatletter\def\@captype{figure}\makeatother
\centering
\begin{minipage}[t]{0.53\textwidth}
\centering
\begin{tikzpicture}[scale=1.15, line cap=round, line join=round]
  \coordinate (v0) at (-2,1.05);
  \coordinate (v1) at (-2.95,-0.65);
  \coordinate (v2) at (-1.05,-0.65);
  \draw[thick,fill=blue!3] (v0)--(v1)--(v2)--cycle;
  \node[above] at (v0) {\(0\)};
  \node[below left] at (v1) {\(1\)};
  \node[below right] at (v2) {\(2\)};
  \node at (-2,-0.12) {\(\sigma\)};
  \node[font=\footnotesize,align=center] at (-2,-1.42)
    {one barycentric\\chamber \(\sigma\simeq\Delta^2\)};
  \draw[-{Stealth[length=2mm]},thick] (-0.25,0.12)--(-0.90,0.12)
    node[midway,above] {\(X\)};
  \begin{scope}[shift={(1,0)}]
    \path[fill=gray!22,draw=black,thick]
      (30:1.40)--(90:1.40)--(150:1.40)--(210:1.40)--
      (270:1.40)--(330:1.40)--cycle;
    \path[fill=white,draw=none]
      (30:0.76)--(90:0.76)--(150:0.76)--(210:0.76)--
      (270:0.76)--(330:0.76)--cycle;
    \node[font=\small] at (60:0.97) {\(\{0\}\)};
    \node[font=\small] at (120:0.97) {\(\{0,1\}\)};
    \node[font=\small] at (180:0.97) {\(\{1\}\)};
    \node[font=\small] at (240:0.97) {\(\{1,2\}\)};
    \node[font=\small] at (300:0.97) {\(\{2\}\)};
    \node[font=\small] at (0:0.87) {\(\{0,2\}\)};
    \node[font=\small] at (0,0) {tile};
  \end{scope}
  \node[font=\small,align=center] at (-0.55,-2.15)
    {(a) A chamber and its resolved tile};
\end{tikzpicture}
\end{minipage}
\hfill
\begin{minipage}[t]{0.43\textwidth}
\centering
\begin{tikzpicture}[scale=1.15, line cap=round, line join=round]
  \fill[blue!12] (0,0) rectangle (1.25,1.25);
  \fill[orange!15] (-1.25,0) rectangle (0,1.25);
  \fill[orange!15] (0,-1.25) rectangle (1.25,0);
  \fill[blue!12] (-1.25,-1.25) rectangle (0,0);
  \fill[gray!25] (-0.10,-1.25) rectangle (0.10,1.25);
  \fill[gray!25] (-1.25,-0.10) rectangle (1.25,0.10);
  \draw[thick] (-1.25,-1.25) rectangle (1.25,1.25);
  \draw[-{Stealth[length=1.6mm]}] (-1.4,0)--(1.5,0)
    node[right,font=\small] {\(Z_S\)};
  \draw[-{Stealth[length=1.6mm]}] (0,-1.4)--(0,1.5)
    node[above,font=\small] {\(Z_T\)};
  \node at (0.68,0.67) {\(\zeta\)};
  \node at (-0.68,0.67) {\(s_S\zeta\)};
  \node at (0.68,-0.67) {\(s_T\zeta\)};
  \node at (-0.68,-0.67) {\(s_Ss_T\zeta\)};
  \node[font=\small] at (1.03,1.04) {\(+\)};
  \node[font=\small] at (-1.03,1.04) {\(-\)};
  \node[font=\small] at (1.03,-1.04) {\(-\)};
  \node[font=\small] at (-1.03,-1.04) {\(+\)};
  \node[font=\small,align=center] at (0,-2.15)
    {(b) \(u_S=Z_S^2,\quad u_T=Z_T^2\)};
\end{tikzpicture}
\end{minipage}
\caption{The reflected resolution in dimension two, following
Gaifullin's permutohedral construction \cite{Gaifullin2008}. In (a), one
barycentric chamber \(\sigma\simeq\Delta^2\) is paired with its resolved
tile. The six facets are indexed by the nonempty proper subsets of
\(\{0,1,2\}\); adjacent labels are nested. The map \(X\) sends the tile
interior diffeomorphically to the chamber interior. In (b), at a corner
with labels \(S\subset T\), a full state \(\zeta=(\sigma,p,t)\) and its three reflected states occupy the four
quadrants of signed coordinates. The commutation
\(s_Ss_T=s_Ts_S\) and the independent bit toggles give four distinct
quadrants of one smooth chart. The displayed signs compare each chamber
sign with that of \(\zeta\): blue agrees and orange is opposite. The gray collar
and strips represent the same seam schematically; their widths are
illustrative, and \(X\) is not affine.}
\label{fig:resolution}
\end{minipage}
\par\medskip


We verify that these are ordinary smooth charts.  On an overlap of two permutations,
a cut \(S_j\) which is not a cut of the other permutation has
\(u_j\geq4^{-2m}\).  To see this, choose a pair of indices on opposite
sides of that cut whose order is inverted in the other permutation.
Their ratio in the first ordering is at most \(4^m u_j\); its inverse
in the second ordering is at most \(4^m\).  The asserted lower bound
follows.  At a cut common to both permutations, the ratio for the
second chart is the first ratio at that cut multiplied by a product
of ratios within the blocks between common cuts.  Those within-block
ratios and their required reciprocals are bounded away from zero on
the overlap.  After taking square roots, the transition at a common
cut is consequently its signed coordinate times a smooth nonvanishing
factor.  All other transition coordinates are smooth monomials with
denominators bounded away from zero.

The sign in this transition is fixed on each region with fixed signs
at the noncommon cuts.  The remaining state relation is a product of
the common reflections, and its product is unique by their free
action on the bits.  Commutation makes this sign rule persist across
the zero coordinates.  Thus the transitions are smooth and invertible
across every corner.  A point incident to \(j\) facets has exactly the
finite star of \(2^j\) states under these reflections, and hence has
one neighborhood of this form rather than several branches.

The signed charts give smooth leaf neighborhoods.  We next check the
topological and transverse properties needed for measured integration.
The leaves are Hausdorff.  Distinct
points with distinct resolved coordinates can be separated in the
tile model.  If their resolved coordinates agree but their finite
stars differ, choose neighborhoods in these two stars avoiding all
nonincident facets; no further identification then joins them.
Each leaf has at most countably many states reachable by the finitely
many moves and at most countably many charts, so it is second countable.
The total space is standard Borel: the quotient of the state--tile
space identifies finite classes at each tile point, for which Borel
representatives can be selected by a fixed enumeration.  More
explicitly, charts can be indexed by the starting chamber, its bit
array, and the flag, with the pairing arrays as transversal parameters.
The bits determine the crossings at nonzero signed coordinates; the
moves can then be inverted to recover the pairing array.  With the
zero convention above this makes the chart maps injective.  Changes
of transversal are the measure-preserving moves on countably many
Borel overlap pieces.  A plaque meets only elementarily many plaques
in this atlas: for each second flag only the states in the two finite
flag reflection groups can occur.  These observations give precisely
the measured flow-box properties of Theorem~\ref{thm:degree}.

The chamber signs orient the leaves.  In a flag chart the Jacobian
sign from positive ratio coordinates to oriented barycentric
coordinates depends only on the permutation.  Changing one signed
coordinate from positive to negative changes the Jacobian sign of
\(u_j=Z_j^2\) once, while the move changes the chamber sign once.
The induced orientation is therefore the same across each facet.
The transverse measures here are nonnegative; the signs enter only
through this orientation in integrals of forms.  The value of
\(\omega\) is constant along each tile leaf, and translation by
\(\Gamma\) reindexes the states and preserves these measures.

\medskip\noindent\emph{The map and the local current.}
Recover the raw \(x_i\) from \eqref{eq:transfer-ratios} by normalizing
the successive products \(1,u_1,u_1u_2,\ldots,u_1\cdots u_m\).
If \(v_{\sigma,i}\) are the vertices of the indicated chamber in its
flat chart, map its tile by
\begin{equation}\label{eq:transfer-tile-map}
 X=\sum_{i=0}^m b_i v_{\sigma,i},\qquad
 b_i=\frac{\exp(-1/x_i)}{\sum_l\exp(-1/x_l)},
\end{equation}
where \(\exp(-1/x_i)\) is extended by zero at \(x_i=0\).  On a facet
with label \(S\), all coefficients outside \(S\) vanish to infinite
order in its signed coordinate.  The vertices in \(S\) are the same
on the paired chambers, and their coefficients are smooth functions
of the squares of the signed coordinates.  Thus \(X\) is smooth
through the gluing.  Its denominator is at least \(e^{-(m+1)}\), since
some \(x_i\geq1/(m+1)\), which also gives elementary upper derivative
bounds in the tile charts.

On the positive simplex the map \(x\mapsto b\) is a diffeomorphism.
Given positive \(b_i\) of sum one, its inverse is obtained from the
unique \(t\in(0,1/\max_i b_i)\) satisfying
\[
                    \sum_i\frac1{-\log(tb_i)}=1;
\]
the left side increases smoothly from zero to infinity.  The inverse
function theorem also gives the smooth inverse.  This diffeomorphism
preserves orientation.  Indeed it extends to the closed simplex
preserving every face, and the linear homotopy to the identity has
the same property, so its relative topological degree is \(+1\).
It follows that the tile interior maps diffeomorphically, with sign
\(\varepsilon_\sigma\), onto the interior of its chamber.  Each
chamber's state probability integrates to one and the tile boundaries
have leafwise measure zero.  Fubini and change of variables therefore
give, for every compactly supported smooth \(m\)-form \(\alpha\),
\[
       \int_{\mathcal A}X^*\alpha
          =\sum_\sigma\varepsilon_\sigma\int_\sigma\alpha.
\]
This proves \eqref{eq:transfer-push}.  To obtain the local identity,
define \(\mathcal A_{\rm orig}\) as the image in
\(X^{-1}(\mathcal T_\delta)\) of tile points belonging to original
chambers.  At an interior point this subset is leaf-open.  At a lower
face, the incidence argument preceding the lemma shows that every state
in its finite reflection star is on the same side.  A neighborhood in
that star can be chosen to keep its image in the strict tube and to
avoid nonincident facets, so it stays on that side.  This proves
leaf-openness and makes the definition unambiguous at glued points;
the finite indexed charts also make it Borel.  Applying the same
change-of-variables calculation only to original tile interiors, for
a form supported in \(\mathcal T_\delta\), gives
\eqref{eq:transfer-local-push}: each original cube is counted with
its coefficient \(D_0u'_C\), and each chamber's state fiber has mass one.

\medskip\noindent\emph{The metric and upper density.}
The signed charts admit a background metric \(k_0\) with elementary
bounds for its coefficients, inverse, and every fixed number of
derivatives in buffered charts.  To construct it, take smooth product
cutoffs supported in the signed boxes \((-2,2)^m\) and equal to one
on \([-1,1]^m\), and average their Euclidean chart metrics with the
normalized cutoff weights.  The smaller boxes cover, by sorted flags.
The overlap bound above bounds the number of terms per plaque.  The
transition formulas and their inverses are products of at most
\(O(m)\) factors, and every required reciprocal at a noncommon cut
is at most \(4^{O(m)}\).  Their fixed-order derivatives consequently
have elementary bounds.  These facts give the claimed bounds for the
average and its inverse, and make the construction Borel and
equivariant.  The same transition and product estimates, applied to
\eqref{eq:transfer-tile-map}, give an elementary upper bound for
\(dX\) and its fixed-order derivatives before any choice of
\(\lambda\).  The affine simplexes have diameter at most \(r\sqrt m\).

Put \(\Theta=\prod_i x_i\).  Choose a smooth \(j=j(\Theta)\) equal
to one for \(\Theta\leq\gamma_0\), zero for
\(\Theta\geq2\gamma_0\), and between zero and one elsewhere, with
the rescaled smooth cutoff estimates.  It is smooth in the signed
charts and the same on both sides of every facet.  Set
\begin{equation}\label{eq:transfer-resolution-metric}
                  k=X^*k_{\rm flat}+j^2 k_0,
       \qquad \text{seam}=\{\Theta\leq2\gamma_0\}.
\end{equation}
We call the complement of this closed seam the straight part.
The metric is positive definite near the tile boundary because \(j=1\) there.  Where
\(j<1\), one has \(\Theta>\gamma_0\), hence \(x_i\geq\gamma_0\) for
every \(i\), and \(X\) is an immersion.  The formula proves
\(\|dX\|_k\leq1\) and equality of metrics off the seam.

The upper density bound requires no inverse estimates.  As quadratic forms,
\(k\leq X^*k_{\rm flat}+k_0\), so its volume density in every sorted
box is bounded by the upper bounds for \(dX,k_0\) already chosen.
Call a common elementary bound \(H_0(q)\).  In a sorted box
\(|Z_j|\leq1\).  If all \(|Z_j|>\lambda\), normalization of the
successive products gives
\[
       \Theta\geq(m+1)^{-(m+1)}
                    \lambda^{2(1+2+\cdots+m)}
              =(m+1)^{-(m+1)}\lambda^{m(m+1)}.
\]
Condition \eqref{eq:transfer-gamma} shows that the seam lies in the
stated union of slabs.  Its coordinate volume in such a box is at
most an elementary factor times \(\lambda\).  The bound \(H_0(q)\)
was fixed before this factor was made small.

The remaining metric estimates may depend on \(\gamma_0\).  On the part
with \(j<1\), all positive ratios are bounded below by powers of
\(\gamma_0\), so the product formulas control the inverse from raw
barycentric to signed coordinates.  The normalized weights in
\eqref{eq:transfer-tile-map} obey
\(b_i\geq e^{-1/\gamma_0}/(m+1)\).  On a tangent vector
\(w=(w_i)\), \(\sum_i w_i=0\), their derivative is
\[
             db_i(w)=b_i\left(\frac{w_i}{x_i^2}
                   -\sum_l b_l\frac{w_l}{x_l^2}\right).
\]
These differences control \(w\): if their common subtracted mean is
\(a\), the relation \(\sum_iw_i=0\) expresses \(a\) as a weighted
average, with weights \(x_i^2\), of the negatives of those differences.
Thus \(|w|\) is bounded by a polynomial in \(m\) times
\(e^{1/\gamma_0}\max_i|db_i(w)|\).  Finally, the affine map from the
barycentric simplex to a simplex in a subdivided cube has elementary
inverse bounds.  Its matrix has \(r\)-scaled rational entries with
bounded denominators, and its nonzero determinant and cofactors give
this by minors.  These estimates control the inverse metric where
\(j<1\); where \(j=1\), \(k_0\) already supplies that control.
Differentiating the formula for \(k\) a fixed number of times now gives
elementary bounds, including the curvature bound, when
\(\gamma_0^{-1}\) is elementary.

\medskip\noindent\emph{Frame interpolation.}
The comparison needs a constant frame ordering tangent and normal axes
on each straight tile, a frame that glues smoothly near facets, and a homotopy to the
identity for the degree calculation.  Choose a Borel skew-symmetric logarithm
\(A_\sigma\) of each \(O_\sigma\) with operator norm at most \(\pi\).
Such a choice follows by decomposing an orthogonal rotation into its
two-dimensional angle blocks with principal angles and, on its
\(-1\)-eigenspace, pairing a Borel choice of orthonormal basis; the
even dimension there follows from determinant one.  Define on that
tile
\[
                   \mathcal R=\exp((1-j)A_\sigma).
\]
It is the identity in a neighborhood of each facet and thus glues
smoothly despite the change of \(\sigma\).  It equals \(O_\sigma\)
where \(j=0\), and replacing \(1-j\) by \(t(1-j)\), \(0\leq t\leq1\),
gives the asserted homotopy.  The cutoff and metric inverse estimates
give the derivative bounds.  The choices can be made on the indexed
chamber data, so they are Borel and equivariant as asserted.
\end{proof}

For our use, choose \(O_\sigma\) to take \(z\)-increments to rotated
grid increments, putting the \(m\) axes tangent to the dual cube first
and the \(n\) complementary axes last.  Use a signed permutation if
needed to have determinant one.  Thus \(O_\sigma=Q_\sigma V\) for a
signed permutation \(Q_\sigma\in SO(q)\) with this order.  Signs will
not affect any row norm or any cubical test.

\subsection{The comparison map and its degree}

The parameter selection below fixes \(\lambda(q)\) and \(\gamma_0(q)\)
deterministically.  For those choices the bounds in
Lemma~\ref{lem:transfer-resolution} give fixed elementary numerical
envelopes \(K_R(q)\) and \(K_{\rm scal}(q)\) such that
\[
 \|d\mathcal R\|_{\infty,k}\leq K_R(q),\qquad
                    (-\Scal_k)_+\leq K_{\rm scal}(q)
\]
for every allowed diagram and every \(V,s,\omega\), rounding parameter,
and state.  These envelopes come from the uniform chart derivative
bounds and do not use \(C_{\rm ent}\).  Fix the single number
\begin{equation}\label{eq:transfer-tau}
 \tau=\tau(q)=
 \frac{1}{300\sqrt q\max\{r^{-1},\delta^{-1},1+K_R(q)\}}.
\end{equation}
It has an elementary reciprocal and is common to all those parameters.
In particular
\(3\tau\sqrt q\leq\min\{r/100,\delta/100,(1+K_R(q))^{-1}\}\).

Fix \(V,s\) and the rounding parameter for the moment.  Take the
product of \(\widetilde M\) and the tile family \(\mathcal A\) over
their shared \(\omega\), and divide by the diagonal \(\Gamma\)-action.
Orient the product with the \(\widetilde M\) directions first and the
\(\mathcal A\) directions second.  The resulting measured leaves have
dimension \(n+m=q\).  Using the leaf-open set
\(\mathcal A_{\rm orig}\) from Lemma~\ref{lem:transfer-resolution},
work only on the open set of pairs \((x,a)\) with
\(a\in\mathcal A_{\rm orig}\) for which \(X(a)\) is in the lift of the ball chart of
\(F(x)\), at distance less than \(\delta/4\).  Define there
\begin{equation}\label{eq:transfer-comparison}
                   G(x,a)=\mathcal R(a)(z_F(x)-z_{X(a)}).
\end{equation}
Here \(z_{X(a)}=z\circ X(a)\).  The requirement on the lifts is part of
the domain; coincident color coordinates on incompatible sheets are
never compared.  The map is invariant under the diagonal action.

Use the band \(\{|G|_\infty<3\tau\}\) in that open domain, and put
\(b_q=\tau(1+3/(4q))\), the inner width in Theorem~\ref{thm:degree}
when \(d=q\).  Let \(\mathcal D\) be its measured degree.

\begin{lemma}[Comparison degree]\label{lem:transfer-degree}
The band in \eqref{eq:transfer-comparison} satisfies the chart and
integration control assumptions of Theorem~\ref{thm:degree}, for each
fixed instance, the product metric \(h+L^2k\) with any fixed \(L>0\),
and the pullback of the entropy weight \(\phi\).  Its degree is Borel
in the auxiliary parameters and satisfies
\begin{equation}\label{eq:transfer-degree}
                         \E_{\rm rd}\mathcal D=D_0I
\end{equation}
for every fixed \(V,s\).
\end{lemma}

\begin{proof}
The local pushforward \eqref{eq:transfer-local-push} will express the
degree using a smoothed form of \(Z_{\rm rd}\).  Its primal-cell periods
are the integer coefficients \(D_0u'_C\).  Their expectations match the periods of
\(D_0\eta\); a compatible piecewise cubical primitive and mollification
then give the same measured pairing with \(F\).  We first verify the
band control that makes these integrations valid.

\emph{The controlled comparison band.}
Before taking the quotient, the product flow boxes use the entropy
transversal \(\omega\) and the
state transversals of \(\mathcal A\).  Variation along the tile keeps
\(\omega\) constant.  These transversals carry the invariant measures
constructed above.  A group element that stabilizes a product leaf
acts freely and properly because it acts so on the \(\widetilde M\)
factor.  Hence the quotient leaves are ordinary oriented Hausdorff
second countable manifolds.  The actions reindex labels and states,
leave the color and flag coordinates unchanged, and preserve the
transverse measures and orientations.

For an inner closed comparison cube there is a uniform margin from
the boundary of the allowed sheet domain.  Indeed \(\mathcal R\) is
orthogonal and \eqref{eq:transfer-tau} gives
\(|z_F-z_X|_2\leq3\tau\sqrt q\ll\delta/4\).  The large witness
color at \(F\) remains present with a uniform small buffer.  In a
flat ball chart a lifted path whose projected coordinates stay in
the ball stays in that same lift, by uniqueness of lifts for a local
diffeomorphism.  The nearby points therefore remain in coherent
sheet charts.  Facets that lead to the mirror are outside this region.
The available smooth bounds for the instance give fixed buffered
coordinate charts for all the data in these regions.

To check the finiteness condition on those charts, choose lifts of a
finite buffered atlas on \(M\).  On their compact closures \(F\) uses
only uniformly finitely many possible labels, by the lifted support
bounds in Theorem~\ref{thm:entropy}.  Any nearby cell lift has its
anchor in a closed one-neighborhood of one of these labels.  The
numbers of projected grid indices, chamber subdivisions, and copies
there are bounded as above.  Thus only finitely many chamber indices
are needed over each such base chart, uniformly for the instance.
Their state probabilities have total mass one per chamber, and the
sum of transverse measures of this covering is finite.  Plaque
overlaps are bounded by the tile overlap bound and the bounded
overlap of the base atlas.  In the quotient only finitely many deck
changes meet two of these buffered base charts.  Subdividing their
inner boxes and, when necessary, restricting to Borel transverse sets
gives the flow boxes stipulated in Theorem~\ref{thm:degree}.  All
metric and data bounds in this verification may depend on the
instance.  It also justifies Fubini and measured Stokes for the
integrations that follow.

\medskip\noindent\emph{The local convolution form.}
Choose a smooth unit-integral top form \(\beta\) on \(\mathbb R^q\)
supported in a Euclidean ball about zero of radius less than
\(\min\{\tau,r/10\}\).  In calculating the degree using \(\beta\),
homotope \(\mathcal R\) to the identity as in
Lemma~\ref{lem:transfer-resolution}.  All maps in this homotopy are
orthogonal on the difference variable.  Their pullbacks of \(\beta\)
are thus supported where the difference has this same small norm,
strictly inside both the sheet domain and the original band.  Measured
Stokes for this supported homotopy shows that it does not change the
degree.

After removing the rotation, fiber integration of
\((z,w)\mapsto z-w\) against \(\beta\) and the original current
\(Z_{\rm rd}\) gives a smooth \(n\)-form \(\alpha_{\rm rd}\) near \(K\).
More explicitly, in a flat chart put \(\mathsf d(z,w)=z-w\).  Integrate
\(\mathsf d^*\beta\) along each compatible \(C^\#\) in the second
variable, multiply by \(D_0u'_C\), and sum.  The sum is locally finite
and uses the original side.  By the local identity
\eqref{eq:transfer-local-push}, Fubini, and change of variables on tile
interiors, the degree is
\begin{equation}\label{eq:transfer-convolution-degree}
                       \mathcal D=\int_{\mathcal M}F^*\alpha_{\rm rd}.
\end{equation}
Boundary cancellation of the dual cubes shows that
\(d\alpha_{\rm rd}=0\) on, for example, the tube of radius
\(\delta/8\) around \(K\).  To verify this directly, apply fiber
Stokes to each cube.  Its boundary terms only involve facets within
the support distance of the point \(z\); these lie in the balanced
tube and cancel with their paired incidence signs.

For every primal \(n\)-cell sufficiently near \(K\) in the same
sheet,
\begin{equation}\label{eq:transfer-convolution-cells}
                         \int_C\alpha_{\rm rd}=D_0u'_C.
\end{equation}
Indeed a primal \(n\)-cell and a complementary dual \(m\)-cell of
this grid meet only when their indices match, and then at their common
center.  If their axis sets differ, a coordinate fixed for both is a
grid vertex for the primal cell and a half-step center for the dual
cell, so their difference in that coordinate is at least \(r/2\).
If their axis sets agree but their grid indices differ, a differing
coordinate gives the same separation.  Such pairs cannot meet the
support of \(\beta\).  For the matching pair, the difference map
from \(C\times C^\#\) contains the support ball in its image with
multiplicity one.  Our fiber integration convention is
\(\int_C\int_{C^\#}\mathsf d^*\beta=
\int_{C\times C^\#}\mathsf d^*\beta\) with the base-first product
orientation.  The differential of \(\mathsf d(z,w)=z-w\) sends this
ordered frame to the primal basis followed by the negative dual basis,
so its orientation is positive by our choice of the dual orientation.
Its integral is therefore
one, proving \eqref{eq:transfer-convolution-cells}.  The small support
and the chart buffers ensure that only the indicated lifts contribute.

\medskip\noindent\emph{The detector pairing.}
Set \(\alpha=\E_{\rm rd}\alpha_{\rm rd}-D_0\eta\).  It is a smooth
closed form on that tube, and \eqref{eq:transfer-rounding},
\eqref{eq:transfer-cell-value}, and
\eqref{eq:transfer-convolution-cells} give \(\int_C\alpha=0\) on all
the primal \(n\)-cells there.  We explain why its measured pairing
with \(F\) vanishes, including when \(F\) meets grid faces.

Use all lifted full \(q\)-cells needed to cover the tube of radius
\(\delta/32\) around \(K\).  Taking \(r\sqrt q\) sufficiently small
as above, these cells lie in the tube of radius \(\delta/8\), all
their \(n\)-faces are selected, and their shared faces fit in
coherent sheet charts.  On one grid interval \([a,a+r]\), project a
function to its linear interpolation between the endpoints and
project a one-form \(f(t)dt\) to the constant one-form with the same
integral.  The homotopy for a one-form is the function
\[
       \mathsf h(f(t)dt)(t)=\int_a^t
              \left(f(v)-r^{-1}\int_a^{a+r}f(s)\,ds\right)dv,
\]
which vanishes at both endpoints; the homotopy on functions is zero.
Direct differentiation gives \(\mathrm{id}-\mathsf P
=d\mathsf h+\mathsf h d\).  Apply these operations one coordinate
at a time, with the ordinary exterior differential signs.  The
tensor-product homotopy still commutes with tangential restriction
to a face, since the one-dimensional primitive vanishes at the
endpoints.  The full projection of an \(n\)-form has, for each
axis set, coefficients obtained by integrating on its \(n\)-faces
and interpolating linearly in the remaining axes.  Hence the full
projection of \(\alpha\) is zero.  Since \(d\alpha=0\), the
homotopy supplies a primitive on each full cube, and its tangential
restrictions agree on shared faces.

This piecewise primitive has bounded coefficients for the instance
and is also a weak primitive on the tube: in distributional Stokes
the terms on each shared face cancel by agreement of its tangential
restrictions.  On a smaller tube mollify it in the flat coordinates.
This can be done with the same small Euclidean convolution in every
chart: the buffered sheets give a unique lift of the convolution ball,
and on overlaps the piecewise primitive and flat coordinates agree.
Its smooth exterior derivatives are the mollifications of the smooth
form \(\alpha\), and converge to \(\alpha\) there.  The construction
is Borel and equivariant and has bounded coefficients for the instance.
Measured Stokes applied to the pullback of each smooth primitive gives
zero; convergence and the uniform chart bounds give
\(\int_{\mathcal M}F^*\alpha=0\).  All expectations and integrals
above are measurable by their locally finite chart formulas, and
their absolute bounds for the instance justify their interchange.
Using \eqref{eq:transfer-form} in
\eqref{eq:transfer-convolution-degree} now proves
\eqref{eq:transfer-degree}.
\end{proof}

\subsection{The two comparison lengths and the final average}

We finish by estimating these degrees.  We will see separately how
the straight part uses the dimension-only Hilbert norm of \(\eta\)
and how the small coordinate volume of the seam pays for its rougher
bound.

\begin{proof}[Proof of Theorem~\ref{thm:scalar-volume}]
The case \(\|M\|=0\) is immediate.  Otherwise we first prove the
conditional target for the data above, choosing \(\Psi_n\) from the
requirements on \(E\) derived below, and then apply it to the entropy
output.  Choose a positive constant \(b_*\), depending only on \(n\),
whose smallness will be fixed in the estimate below, and put
\begin{equation}\label{eq:transfer-lengths}
       L_T=b_*\sqrt E,\qquad L_N=b_*\sqrt{qE},\qquad
                    h_{\rm tot}=h+L_T^2k.
\end{equation}
Once the degree hypotheses are verified, its length denominator, the
product volume, and the entropy contribution will combine as
\begin{equation}\label{eq:transfer-prefactor}
 \frac{L_T^mE^{n/2}}{L_T^mL_N^n}\,e^{-W}dV_h
                         =b_*^{-n}q^{-n/2}\rho.
\end{equation}
This is the cancellation of the tangent volume and all powers of
\(E\).  The estimates below justify the degree application and bound
the straight angular sum by a dimensional multiple of \(q^{n/2}\),
while the collar volume controls the seam.

Use the weight \(\phi\), constant in the tile directions.  The
product scalar formula and the definition of \(\wscal\) give
\(\wscal(h_{\rm tot},\phi)=\wscal(h,\phi)+L_T^{-2}\Scal_k\).
By the definition of \(K_{\rm scal}(q)\), the lower bound
\(E\geq b_*^{-2}K_{\rm scal}(q)\) ensures
\(L_T^{-2}\Scal_k\geq-1\).  By \eqref{eq:transfer-entropy},
\begin{equation}\label{eq:transfer-total-scalar}
 \wscal(h_{\rm tot},\phi)\geq-C_{\rm ent}-1+2W-n\log E+E\mathfrak e.
\end{equation}

Assign \(L_i=L_T\) to the first \(m\) comparison rows and
\(L_i=L_N\) to the last \(n\) rows.  On the band, the derivative in
tile directions of each row of \(G\), measured with \(k\), has norm
at most two.  The \(dX\) term has norm at most one, and the
\(d\mathcal R\) term has norm at most
\(3\tau\sqrt q\|d\mathcal R\|_k\leq1\) by
\eqref{eq:transfer-tau}.  The horizontal derivative is a row of a
rotation of \(dz_F\).  It follows that
\[
 L_i|dG_i|_{h_{\rm tot}}
       \leq C_n\bigl(\sqrt q+\sqrt{qE}\,E^{-1/100}\bigr).
\]
Take \(B\) to be this upper bound increased by two.  Then
\(L_i|dG_i|\leq B-1=B-q/q\).  The horizon condition of
Theorem~\ref{thm:degree} holds after an elementary increase of \(E\):
it suffices to test the smallest length \(L_T\), and, for \(E\geq1\),
\[
       \frac{C(1+B)^{2/3}}{L_T}
          \leq c_{{\rm hor},n}q^{1/3}E^{-13/75}
          <\frac{\tau}{10q}.
\]
Here \(c_{{\rm hor},n}\) depends only on the absolute horizon constant
and the dimensional choice of \(b_*\), and
\(13/75=1/2-(2/3)(1/2-1/100)>0\).  Thus the displayed strict horizon
inequality follows if
\(E\geq(20c_{{\rm hor},n}q^{4/3}/\tau)^{75/13}\).
We may therefore apply Theorem~\ref{thm:degree}
with \(d=q\).

Write \(H_i=2(L_i|dG_i|-1)_+\), as in that theorem.  At a straight
point in a tile, let \(\mathcal I\) be the \(n\)-set of grid axes
normal to its dual cube.  Put \(A_i=(Vdz_F)_i\), with covector norms
measured by \(h\), and define
\begin{equation}\label{eq:transfer-U}
       U_{\mathcal I}^2=
        \frac q{\mathfrak e}\sum_{i\in\mathcal I}|A_i|_h^2
        \quad\text{if }\mathfrak e>0,
        \qquad U_{\mathcal I}=0\quad\text{if }\mathfrak e=0.
\end{equation}
Here \(\sum_i|A_i|^2=\mathfrak e\).  At a straight point
\(k=X^*k_{\rm flat}\) and \(\mathcal R=Q_\sigma V\) is constant.
For a tangent comparison row corresponding to axis \(i\), the squared
row norm is exactly \(L_T^{-2}+|A_i|^2\).  Thus
\[
  \sum_{\rm tangent}H_i
      \leq L_T^2\mathfrak e=b_*^2E\mathfrak e,
\]
by \(2(\sqrt{1+t}-1)\leq t\).  Normal rows have no tile derivative,
so Cauchy--Schwarz gives
\[
  \sum_{\rm normal}H_i
    \leq2L_N\sum_{i\in\mathcal I}|A_i|
    \leq2b_*\sqrt n\sqrt{E\mathfrak e}\,U_{\mathcal I}.
\]
Completing the square in \(\sqrt{E\mathfrak e}\), for \(b_*<1\),
shows that
\begin{equation}\label{eq:transfer-straight-penalty}
       \frac{\sum_iH_i-E\mathfrak e}{2}
          \leq c_n U_{\mathcal I}^2,
       \qquad c_n=\frac{n b_*^2}{2(1-b_*^2)}.
\end{equation}

Choose \(b_*\) sufficiently small that
\begin{equation}\label{eq:transfer-sphere}
                  \E_V\exp(2c_nU_{\mathcal I}^2)\leq C_n
\end{equation}
for every fixed horizontal differential and every fixed
\(\mathcal I\).  Here is a direct verification of this choice.  If
\(u\) is one coordinate of a uniform point on \(S^{q-1}\), its sphere
moments satisfy
\[
 \E(q u^2)^j
   =\frac{q^j(1\cdot3\cdots(2j-1))}{q(q+2)\cdots(q+2j-2)}
   \leq2^j j!.
\]
Consequently \(\E e^{tqu^2}\leq(1-2t)^{-1}\) for \(0\leq t<1/2\).
The singular-value decomposition of \(dz_F\) expresses each
\(q|A_i|^2/\mathfrak e\) as a convex combination of such random
variables, with its singular-value weights.  Jensen's inequality
gives the same exponential bound for that convex combination.
Another application of Jensen to the sum of the \(n\) rows in
\eqref{eq:transfer-U} bounds its exponential moment by
\((1-2nt)^{-1}\).  Choosing, for instance, \(2nc_n<1/4\) proves
\eqref{eq:transfer-sphere}.  No independence of the rows of \(V\)
is used.

At a seam point the uniform tile row bound gives a rougher estimate
which is still independent of the instance.  By the triangle inequality
for each row, \((a-1)_+\leq a\), and Cauchy--Schwarz on the horizontal
rows,
\[
 \sum_iH_i\leq 2L_N\sqrt{q\mathfrak e}
                   +4q L_N/L_T
       \leq2b_*q\sqrt{E\mathfrak e}+4q^{3/2}.
\]
Completing the square once more gives
\begin{equation}\label{eq:transfer-seam-penalty}
 \frac{\sum_iH_i-E\mathfrak e}{2}\leq c_{{\rm seam},n}q^2,
                 \qquad c_{{\rm seam},n}=2+\frac{b_*^2}{2}.
\end{equation}

Apply the degree inequality, then use
\eqref{eq:transfer-total-scalar}.  Its loss is \(\exp(Cq/q)\), an
absolute constant.  The product volume form is
\(dV_{h_{\rm tot}}=L_T^m dV_h\,dV_k\), so the measure factor is
exactly \eqref{eq:transfer-prefactor}.  Together with
\eqref{eq:transfer-degree} and \(|\E\mathcal D|\leq\E|\mathcal D|\),
this gives the following inequality, with
\(C_{\rm out}=\exp(C_{\rm deg}+(C_{\rm ent}+1)/2)b_*^{-n}\), where
\(C_{\rm deg}\) is the absolute exponent in Theorem~\ref{thm:degree}:
\begin{equation}\label{eq:transfer-average}
 \begin{split}
 D_0I(2\tau)^q\leq C_{\rm out}q^{-n/2}
  \int_{\mathcal M}\rho\,
   \E_{V,s,\rm rd}
     \int_{\{a:\ |G(x,a)|_\infty<b_q\}}
                       \mathcal B(x,a)\,dV_k(a),
 \end{split}
\end{equation}
where the fiber integral includes its state transversals and only the
compatible original-side band, and
\[
 \mathcal B(x,a)=
 \begin{cases}
    \exp(c_nU_{\mathcal I}^2),&a\text{ straight over axes }\mathcal I,\\
    \exp(c_{{\rm seam},n}q^2),&a\text{ in the seam}.
 \end{cases}
\]
The entropy constant enters only through this final prefactor, not the
geometric bounds used below.  The length and volume prefactor has no
exponential loss in \(m\); the seam exponential will be absorbed by
the choice of \(\lambda\).

Fix \(x,\omega\), and put \(y=F(x)\).  At a coupled pair in
\eqref{eq:transfer-average}, a fixed projected grid index has at most
one possible cell lift in the ball chart of \(y\).  Its full primal
and dual cells fit in that chart by \eqref{eq:transfer-tau} and the
buffers, and the lift there is unique.  Moreover each point of its
primal cell is within \(10r\sqrt q\) of \(y\): the dual point is
within \(b_q\sqrt q\) of \(y\), and the dual and primal extents are
at most \(r\sqrt m/2\) and \(r\sqrt n/2\).  The variation choice
for \(r\) and the rounding error now give the copy-count estimate
\begin{equation}\label{eq:transfer-copies}
       D_0|u'_C|\leq D_0r^n
                        \bigl(|\eta_{\mathcal I}(y)|+2q^{-n}\bigr),
\end{equation}
where the component is in the \(V\)-rotated frame.

On a straight tile the \(k\)-volume is exactly the flat volume of
its image in the chamber.  The signed permutation in \(G\) does not
change its cube test, which is
\(|\widehat z_y-\widehat z_X|_\infty<b_q\).  Summing the chamber
pieces in a copy is thus bounded by the \(m\)-area of its whole dual
cube inside that test; the state probabilities for each chamber
integrate to one.  For a fixed axis set \(\mathcal I\), summing over
all projected centers and averaging the shift gives exactly
\begin{equation}\label{eq:transfer-shift}
  r^{-q}\int_{\mathbb R^q}\int_{[-r/2,r/2]^m}
       \indic_{\{|\widehat z_y-c-t|_\infty<b_q\}}\,dt\,dc
                    =r^{-q}r^m(2b_q)^q.
\end{equation}
Here \(t\) is placed in the complementary axes.  This identity is
just the density \(r^{-q}\) of centers under uniform shift; extending
the sum to the full grid is an upper bound for the selected lifted
cells.

The form bound holds in every rotated frame.  Using Cauchy--Schwarz
jointly in \(V\) and the \(\binom qn\) axis sets, and then
\eqref{eq:transfer-sphere}, we obtain
\begin{equation}\label{eq:transfer-form-average}
 \begin{split}
 \E_V\sum_{|\mathcal I|=n}
   (|\eta_{\mathcal I}(y)|+2q^{-n})e^{c_nU_{\mathcal I}^2}
 &\leq
   \left(\E_V\sum_{\mathcal I}|\eta_{\mathcal I}(y)|^2\right)^{1/2}
   \left(\E_V\sum_{\mathcal I}e^{2c_nU_{\mathcal I}^2}\right)^{1/2}
       +C_nq^{-n}\binom qn\\
 &\leq C_nq^{n/2}.
 \end{split}
\end{equation}
This use of Cauchy--Schwarz allows the form components and row norms
to depend on the same rotation.  Combining
\eqref{eq:transfer-copies}--\eqref{eq:transfer-form-average}, the
straight part of the inner expectation in
\eqref{eq:transfer-average} is at most
\(C_nD_0q^{n/2}(2b_q)^q\), because \(r^{n+m-q}=1\).  Multiplying
by its prefactor \(q^{-n/2}\) bounds that part, for a fixed
dimension-only constant \(C_{{\rm str},n}\), by
\begin{equation}\label{eq:transfer-straight-average}
             C_{{\rm str},n}D_0(2\tau)^q,
\end{equation}
since \((b_q/\tau)^q=(1+3/(4q))^q\leq e^{3/4}\).

For the seam use only the consequence
\(|z_X-z_y|_2\leq b_q\sqrt q\) of the cube test.  Discard the tile
boundaries, which have leafwise measure zero.  The remaining set is
the disjoint union of tile interiors indexed by their \emph{actual}
chamber \(\sigma\) and its state arrays.  Partition each such interior
into positive sorted-coordinate pieces, for example by taking the
first sorted permutation in a fixed order.  On these pieces
\(0<Z_j\leq1\), so no signed-chart reflection changes the chamber.
The conditional state measure over each actual chamber has mass one,
and its \(k\)-density is at most the elementary \(H_0(q)\) chosen
before \(\lambda\).

Retain the compatible-sheet restriction at this point.  For a fixed
projected grid index the uniqueness proved before
\eqref{eq:transfer-copies} allows at most one compatible geometric
lift.  Apply the rough copy bound
\(D_0|u'_C|\leq C_{\eta,n}D_0r^n\) from that equation, where
\(C_{\eta,n}\) depends only on the form norm, and use the uniform
density bound.  Only after these bounds may we drop the restrictions
on available sheets, roundings, and chambers and sum over all projected
centers and chamber shapes.

For a fixed actual chamber shape, normal axis set \(\mathcal I\), and
positive sorted flag, \eqref{eq:transfer-tile-map} has rotated position
\[
 \widehat z_X=s+r\left(k+\tfrac12\indic_{\mathcal I}\right)
                       +d_{\mathcal I,\mathrm{shape},\mathrm{flag}}(Z).
\]
Here \(\indic_{\mathcal I}\) is the coordinate indicator vector of
\(\mathcal I\).
The displacement \(d\) is independent of the pairing and bit arrays,
of \(\lambda,\gamma_0\), and of the shift; it is determined by the
fixed chamber shape and the barycentric formula.  For each fixed \(Z\)
the shift integral of the rough test is consequently
\[
 \begin{split}
 \frac1{r^q}\int_{[0,r)^q}\sum_{k\in\mathbb Z^q}
  \indic_{\{|s+r(k+\frac12\indic_{\mathcal I})+d(Z)-\widehat z_y|_2
                  \leq b_q\sqrt q\}}\,ds
 &=r^{-q}\Vol_{\mathbb R^q}(B_{b_q\sqrt q})\\
 &\leq r^{-q}(4\tau\sqrt q)^q.
 \end{split}
\]
This calculation is valid despite the shift dependence of the actual
diagram, because the copy, lift, and density bounds were imposed before
the enlargement to all centers.  Tonelli's theorem applies to these
nonnegative bounds.

In each positive sorted piece the seam lies in
\(\bigcup_j\{Z_j\leq\lambda\}\), of coordinate volume at most
\(m\lambda\).  A cube has \(2^m m!\) barycentric top chambers and
at most \((m+1)!\) sorted flags per chamber.  Thus the seam part of
the inner expectation is bounded by
\begin{equation}\label{eq:transfer-seam-average}
                    C_*(q)\lambda D_0(2\tau)^q,
\end{equation}
where one may take the elementary envelope
\[
 C_*(q)=\max\left\{1,\,
 C_{\eta,n}H_0(q)\binom qn(2^m m!)(m+1)!\,m\,
 r^{n-q}(2\sqrt q)^q\exp(c_{{\rm seam},n}q^2)\right\}.
\]
Every factor here was bounded before choosing \(\lambda\).  Thus
\(C_*\) is independent of \(\lambda,\gamma_0,\tau\), of the bounds
\(K_R,K_{\rm scal}\), and of \(C_{\rm ent}\).  The powers \(\tau^q\) in the shift
count have already canceled those in \((2\tau)^q\).

The two fiber estimates now have the same normalization.  Insert
\eqref{eq:transfer-straight-average} and
\eqref{eq:transfer-seam-average} into \eqref{eq:transfer-average},
using \(q^{-n/2}\leq1\) for the seam, and cancel the positive factor
\(D_0(2\tau)^q\).  For the conditional data and parameter bounds used
so far, this gives
\begin{equation}\label{eq:transfer-before-choices}
       I\leq C_{\rm out}\bigl(C_{{\rm str},n}+C_*(q)\lambda\bigr)
                           \int_{\mathcal M}\rho.
\end{equation}

We now fix these bounds before invoking the entropy theorem.  For each
integer \(q\geq n+1\), the form derivative bound first gives a
deterministic \(r(q)\) with elementary reciprocal.  The size bound for
the universal type list, \(D_0\), incidence counts, and the bounds for
\(k_0,dX\), including \(H_0(q)\), have elementary envelopes uniform
over all permitted graphs and cocycles.  Fix \(b_*\) using
\eqref{eq:transfer-sphere}; this determines the dimensional constants
\(c_{{\rm seam},n}\) and \(c_{{\rm hor},n}\), independently of
\(C_{\rm ent}\).  The displayed formula for \(C_*(q)\) is consequently
available at this stage.  Set
\[
 \lambda(q)=\frac1{4C_*(q)},\qquad
 \gamma_0(q)=\frac14(m+1)^{-(m+1)}\lambda(q)^{m(m+1)}.
\]
These deterministic choices have elementary reciprocals and satisfy
\eqref{eq:transfer-gamma}.  The resolution estimates now give common
elementary envelopes for its fixed-order chart derivatives, hence
\(K_R(q),K_{\rm scal}(q)\), and then the common \(\tau(q)\) in
\eqref{eq:transfer-tau}.  None of these choices uses the entropy constant.

For integer \(j\geq n+1\), define
\[
 \Psi_n^\circ(j)=\max\left\{1,\ b_*^{-2}K_{\rm scal}(j),
       \left(\frac{20c_{{\rm hor},n}j^{4/3}}{\tau(j)}\right)^{75/13}
                      \right\},
\]
and define a majorant on all \(t\in[1,\infty)\) by
\[
 \Psi_n(t)=1+\max_{\substack{j\in\mathbb Z\\
                   n+1\leq j\leq\max\{n+1,\lceil t\rceil\}}}
                            \Psi_n^\circ(j).
\]
Every expression used here has an elementary upper bound; the finite
maximum is bounded by an increasing elementary envelope at
\(\max\{n+1,\lceil t\rceil\}\).  Thus \(\Psi_n\) is an elementary
majorant fixed solely from \(n\), and it dominates both the curvature
and strict horizon requirements.  Hence every conditional datum with
\(E\geq\Psi_n(q)\) satisfies all hypotheses used in
\eqref{eq:transfer-before-choices}.  Our choice
\(C_*(q)\lambda(q)=1/4\) and the definition of \(C_{\rm out}\) prove
\eqref{eq:transfer-conditional-target} with
\[
       K_n=e^{C_{\rm deg}+1/2}b_*^{-n}
                             \bigl(C_{{\rm str},n}+1/4\bigr).
\]
This closes the conditional target for all the specified data before
any entropy family is chosen.

Now apply Theorem~\ref{thm:entropy} with \(\mathcal B=\Psi_n\),
\(E_{\min}=1\), the detector root map as \(F^{\mathrm{in}}\), and the
one-point transverse space \(\Omega_0\).  The detector construction uses
finitely many compactly supported bump representatives and their deck
translates, giving the required uniform support and derivative bounds in
a lifted atlas and the equivariant compact support sets.  After pulling
the detector map to \(\Omega\) and embedding its labels equivariantly
among the output's finite copies, the measured affine probability homotopy
of Section~2 preserves \(I\); the output therefore supplies the conditional
data with \(E>\Psi_n(q)\).  Its constant
\(C_{\rm ent}=C_{\rm ent}(n,\Psi_n)\) depends only on \(n\).  Since
\(\rho\leq dV_g\) and the transverse measure over \(M\) is a
probability measure,
\[
           I\leq K_ne^{C_{\rm ent}/2}\int_{\mathcal M}\rho
             \leq K_ne^{C_{\rm ent}/2}\Vol_g(M).
\]
Finally \(I\geq\|M\|/2\), so the constant
\(C_n=2K_ne^{C_{\rm ent}/2}\) proves the theorem.
\end{proof}

\begin{proof}[Proof of Theorem~\ref{thm:main}]
Fix \(n\geq3\), and let \(C_n>0\) be the constant for the normalized
scalar-curvature estimate in Theorem~\ref{thm:scalar-volume}.
Proposition~\ref{prop:yamabe-reduction} gives
\[
              \|M\|\leq C_n\int_M(\Scal_g^-)^{n/2}\,dV_g
\]
for every metric in the theorem.  Taking \(a_n=C_n^{-1}>0\) proves
the stated integral bound.
\end{proof}
